% AI-produced by GPT-6 Astra (OpenAI, ChatGPT Pro mode), 2026-11-19.
% Fragment for the root paper. Environments and natbib are supplied there.

\section{B.3. Finite figures and event structures: the missing distributivity hypothesis}

\begin{definition}[Known: the two irreducibility conditions]\label{SA21-BM-def-prime}
Let \((\mathcal F,\leq)\) be a nonempty partially ordered set with least
element \(0\) and suprema of bounded families. An element \(m\ne0\) is
\emph{completely join-irreducible} if
\[
 m=\bigvee A\quad\Longrightarrow\quad m\in A
\]
whenever the supremum exists. It is \emph{prime for bounded joins} if
\[
 m\leq\bigvee A\quad\Longrightarrow\quad
 (\exists a\in A)\ m\leq a
\]
for every bounded family. The definition in 156-4, scan 54, is the
first condition, with equality on the left. Write \(M\) for these
irreducible figures and \(\delta(F)=\{m\in M:m\leq F\}\) for deployment.
Here a \emph{finite figure system} means that each \(\delta(F)\) is
finite, not that \(M\) or \(\mathcal F\) is finite. We retain C1--C3:
bounded suprema exist, \(F=\bigvee\delta(F)\), and two figures have
a common upper bound exactly when every pair of their multistrates
has a common upper bound.
\end{definition}

\begin{counterexample}[Refuted: R12 and the literal form of L1]\label{SA21-BM-r12}
The literal finite-figure axioms C1--C3 do not imply representation
by finite downward-closed compatibility cliques. The failure persists
after adding C4, C4 bis, and the C5--C8 subdivision axioms, including
C8\(_0\), as stated in 156-4, scans 65--87.
\end{counterexample}

\begin{proof}
Take the six-element lattice
\[
 \mathcal F=\{0,d,a,b,c,1\},\qquad 0<d<a,b,c<1,
\]
where \(a,b,c\) are pairwise incomparable. All omitted comparisons
follow by reflexivity or transitivity. Distinct members of
\(\{a,b,c\}\) have join \(1\) and meet \(d\). All families have
suprema, so C1 holds. The equality-based definition gives
\[
 M=\{d,a,b,c\},\quad
 \delta(0)=\varnothing,\quad\delta(d)=\{d\},\quad
 \delta(u)=\{d,u\}\ (u=a,b,c),\quad\delta(1)=M.
\]
These sets generate their figures, proving C2. Every pair of figures
has upper bound \(1\); compatibility is universal and C3 holds.
The incidence poset \(M\) is connected, since \(d<a,b,c\).
Consequently C4 is vacuous and C4 bis is the identity decomposition
into one component.

Define subdivision by equality. The refinement induced by taking
a subfigure of a subdivision is then \(\leq\). C5, invariance of
disjointness under subdivision, is immediate. In C6, the subdivision
of \(G\leq F\) induced from \(F'=F\) is uniquely \(G'=G\).
Induction is the identity on the principal ideal of \(F\), so
preserves every bounded supremum and satisfies the multistrate
formula. For C7\(_0\), \(G'=G\) extends by \(F'=F\); the stronger
canonical choice is equally forced. A compatible family of
subdivisions in C8 consists of the original cover members, whose
given supremum is its own subdivision. Both zigzag conditions of
C8\(_0\) hold because equality in either subdivision position
forces equality throughout. Moderation holds since each figure
has only its identity subdivision. These verifications agree
with the observation in 156-4, scan 81, that C5--C7 are
tautological for equality subdivisions.

Nevertheless \(\{d,a,b\}\) is a finite downward-closed compatibility
clique in \(M\) and is not a deployment:
\[
 \delta(a\vee b)=\{d,a,b,c\}
 \ne\delta(a)\cup\delta(b)=\{d,a,b\}.
\]
The assertion in 156-4, scan 55, that deployment sends suprema
to unions is therefore false under its preceding hypotheses.
More intrinsically, this lattice is not the finite configuration
poset of any prime event structure. The ideal below a
configuration is the downset lattice of its causal poset and
is distributive, whereas here
\[
 a\wedge(b\vee c)=a,\qquad
 (a\wedge b)\vee(a\wedge c)=d.
\]
The program \texttt{bmodels\_checks.py} checks all families in
this specified lattice and the displayed axioms. This is a
counterexample to the C1--C8 package; it does not satisfy the
later requirement of a lieu in every open shadow.
\end{proof}

\begin{theorem}[Proved: exact corrected event-structure representation]\label{SA21-BM-event-iff}
For a finite-figure system satisfying C1--C3, the following are equivalent:
\begin{enumerate}
\item Every principal ideal \(\mathcal F_{\leq F}\) is distributive.
\item Every multistrate is prime for bounded joins.
\item Deployment is an isomorphism of ordered sets
\[
 \mathcal F\cong
 \{A\subseteq M:A\text{ is finite, downward closed, and an }R
   \text{-clique}\},
\]
where \(mRn\) means that \(m,n\) have a common upper bound.
\item \(\mathcal F\) is the finite configuration poset of a
binary-conflict prime event structure with finite causes.
\end{enumerate}
The event structure is recovered, uniquely up to a bijection
preserving causality and conflict, by
\[
 E=M,\qquad \leq_E=\leq|_M,\qquad
 m\mathbin{\#}n\ \Longleftrightarrow\ \neg(mRn).
\]
These conditions identify the whole representable subclass of
literal finite contrées. Adding C4 corresponds exactly to
requiring all conflict to occur within undirected causality
components. Adding C4 bis while retaining finite figures
requires that there be only finitely many nonempty such components.
\end{theorem}

\begin{proof}
C2 makes deployment order reflecting: if
\(\delta(G)\subseteq\delta(F)\), then
\(G=\bigvee\delta(G)\leq F\). Every principal ideal is therefore
finite, since it injects into \(\mathcal P(\delta(F))\).
C1 makes it a lattice; a meet is the supremum of its common
lower bounds.

Assume (1). For \(m\leq G\vee H\), distributivity in a common
bounding ideal gives
\[
 m=(m\wedge G)\vee(m\wedge H).
\]
Irreducibility gives \(m\leq G\) or \(m\leq H\). Induct on a
finite family; any bounded family takes values in a finite
ideal, so this also proves the arbitrary bounded version.
The empty family causes no exception since \(m\ne0\).
This is the local bounded form of
\citet[Chapter IX, \S4, Lemma 1 and Theorem 5, p.~139]{Birkhoff1948}.

Assume (2). Then
\(\delta(\bigvee_iF_i)=\bigcup_i\delta(F_i)\) for bounded
families. Deployments are finite downward-closed cliques.
Conversely let \(A\subseteq M\) be such a clique. Construct
its join successively. Once the join of \(a_1,\ldots,a_k\)
exists, its deployment is
\(\bigcup_{i\leq k}\delta(a_i)\). Every member of this union
is compatible with \(a_{k+1}\), by downward heredity of
compatibility. C3 makes the partial join compatible with
\(a_{k+1}\), so C1 supplies their join. Thus \(\bigvee A\)
exists, and primality and downward closure give
\[
 \delta(\bigvee A)=\bigcup_{a\in A}\delta(a)=A.
\]
Use \(0\) if \(A=\varnothing\). This proves (3).
Under (3), the ideal below \(F\) is the lattice of all downsets
of the finite poset \(\delta(F)\); meets and joins are
intersection and union, proving (1).

The relation \(R\) is reflexive, symmetric and downward
hereditary. Its complement \(\#\) is therefore symmetric,
irreflexive and upward hereditary. Each principal causal
ideal is \(\delta(m)\), hence finite and conflict free.
Thus (3) is exactly the finite configuration definition
of \citet[Definition 14, pp.~99--100]{NPW1981}, with finite
causes additionally imposed.

Conversely, for a binary event structure with finite causes,
the finite configurations are the finite downward-closed
conflict-free sets. A family bounded by a finite configuration
has union as its supremum. Each configuration is the union
of its principal causes \(\downarrow e\). These are precisely
the nonzero completely join-irreducible configurations:
membership of \(e\) proves their prime property, and
a configuration not equal to any of its principal causes
is their proper join. Compatibility is absence of cross
conflict, so C3 holds. Principal ideals are downset lattices
and are distributive. The maps \(e\mapsto\downarrow e\)
and \(F\mapsto\delta(F)\) make the reconstructions inverse
up to the stated isomorphisms. We identify objects and
isomorphisms here, not an unspecified category of morphisms.

Incidence components of \(M\) are causal components. C4
says precisely that events from different components
never conflict. Configurations then decompose uniquely
as component configurations with finite support.
If there are finitely many components, this is the
whole product, giving C4 bis. Conversely, choosing a
nonempty principal cause in each of infinitely many
components yields a tuple with infinite support, which
is not the trace of a finite configuration. When \(E\)
is empty the empty product is the one-element
configuration poset, as required.
\end{proof}

\begin{counterexample}[Refuted: omission of the word binary]\label{SA21-BM-binary}
Even after correcting R12 by local distributivity, the
corresponding class is the binary-conflict class, not all
prime event structures under the finite-consistency convention.
\end{counterexample}

\begin{proof}
\citet[Definition 1.3.4, pp.~335--336]{Winskel1987} allows
a finite consistency predicate \(\mathrm{Con}\).
Take three events \(e_0,e_1,e_2\), equality causality, and
\[
 \mathrm{Con}=\{A\subseteq\{e_0,e_1,e_2\}:|A|\leq2\}.
\]
Singleton consistency, downward closure, closure under
adding causes, and finite causes all hold. All pairs
are consistent but the triple is not. Hence the
configurations \(\{e_0,e_1\}\), \(\{e_2\}\) violate C3.
No binary relation represents this consistency
predicate, since consistency of all pairs would force
empty conflict. The program checks these axioms and
all eight undirected conflict relations on the three
events. The precise subclass is that of
\citet[\S3.4.6, p.~378]{Winskel1987}.
\end{proof}

\begin{proposition}[Proved: configuration domains and supports]\label{SA21-BM-domain}
For the corrected class, let \(\mathcal C_f(E)\) denote finite
configurations and \(\mathcal C(E)\) all downward-closed
conflict-free subsets. Then
\[
 \mathcal C(E)\cong\operatorname{Idl}(\mathcal C_f(E)),
\]
where ideals are nonempty directed downward-closed sets.
The compact elements are the finite configurations,
and the complete primes are \(\downarrow e\).
This is the coherent prime-algebraic domain correspondence
of \citet[Theorems 8--9, Definitions 17--18, p.~102]{NPW1981}.
For the whole-figure disjointness of 156-4 set
\[
 e\perp_E f\quad\Longleftrightarrow\quad
 \neg(e\#f)\text{ and }\downarrow e\cap\downarrow f=\varnothing.
\]
The support ortholattice of finite configurations is
canonically isomorphic to the closed double polars
of \((E,\perp_E)\), and to the support ortholattice
using all configurations. This concerns whole-figure
disjointness, not the later primitive internal disjointness.
If the empty configuration is retained as an orthogonality
vertex, it is universally self-orthogonal, and every closed
figure polar contains it. Removing this one radical vertex
gives the ordinary irreflexive orthoset convention and the
same support ortholattice.
\end{proposition}

\begin{proof}
Finite configurations contained in a configuration \(x\)
form an ideal: finite unions remain configurations.
Their union is \(x\), because every \(e\in x\) lies
in \(\downarrow e\subseteq x\), which is finite.
Conversely, the union of a directed ideal is downward
closed, and any two of its events lie together in
one member and so do not conflict. Each finite
subconfiguration of this union lies in a member
by finite directedness. These operations are
inverse order isomorphisms.

Directed suprema are unions. Finite configurations
are compact by the finite-witness argument.
An infinite configuration is the directed union
of its finite subconfigurations, so cannot be
compact. Membership of \(e\) proves complete
primality of \(\downarrow e\); decomposition
into these principals proves there are no
others. A pairwise compatible family of
configurations has conflict-free union, proving
coherence and checking the cited theorem's
hypotheses in this setting.

For configurations \(x,y\),
\[
 x\perp y
 \ \Longleftrightarrow\
 x\cup y\in\mathcal C(E),\ x\cap y=\varnothing
 \ \Longleftrightarrow\
 (\forall e\in x)(\forall f\in y)\ e\perp_Ef .
\]
The forward implication uses causal closure.
Conversely cross nonconflict gives a configuration
union; an event in \(x\cap y\) would imply
\(e\perp_Ee\), impossible.

For \(A\subseteq E\) write \(A^{\perp_E}\) for
its event polar and set
\(K_A=\{x\in\mathcal C_f(E):x\subseteq A\}\).
Event polars are downward closed, since
\(\perp_E\) is downward hereditary.
If \(A=A^{\perp_E\perp_E}\), every \(e\in A\)
occurs in \(\downarrow e\in K_A\), whence
\[
 K_A^\perp=K_{A^{\perp_E}}.
\]
For any configuration family \(\mathcal U\),
the pairwise criterion gives
\[
 \mathcal U^\perp=K_{(\bigcup\mathcal U)^{\perp_E}}.
\]
Thus \(A\mapsto K_A\), with inverse union,
is a bijection between closed event polars
and closed figure polars. It preserves
inclusion and orthocomplement and hence
all complete ortholattice operations.
The same argument with all configurations
proves the final assertion. In the convention retaining
the empty figure, \(K_\varnothing=\{\varnothing\}\)
is the bottom closed polar. Removing the empty figure
from every polar is an order-and-complement isomorphism
to the convention with bottom the empty family.
\end{proof}

\begin{proposition}[Proved: refinement is additional domain structure]\label{SA21-BM-refinement}
The event structure determines incidence, compatibility
and the preceding whole-figure supports, but not refinement,
subdivision, lieux, or primitive internal disjointness.
Given multistrate refinement satisfying At 4, its translation is
\[
 g\ll f\quad\Longleftrightarrow\quad
 (\forall e\in g)(\exists d\in f)\ e\ll d .
\]
Under At 13, subdivision translates as
\[
 g\preceq f\quad\Longleftrightarrow\quad
 g\ll f\text{ and }g^\perp=f^\perp,
\]
equivalently maximality for configuration inclusion among
refinements of \(f\).
\end{proposition}

\begin{proof}
The first formula is At 4 through deployment.
Equality of supports is equivalent to equality of
annihilator families by triple-polar equality.
At 13 gives the maximality equivalence of 156-6,
scans 84--85, proved in Part B.1.

For nonuniqueness, take two events \(p,q\) with
equality causality and conflict \(p\#q\).
Configurations are \(\varnothing,\{p\},\{q\}\).
On \(M=\{p,q\}\), let incidence be equality
and primitive internal disjointness empty.
Use equality refinement in the first system;
in the second also put \(p\ll q\).
Both are partial orders containing incidence.
The target's incidence ideal is a singleton,
so the least carrier exists and strong
refinement is the same relation. Mag 3 is
vacuous, and Mag 4 is vacuous on distinct
incident substrata. The first system has
lieux \(p,q\); the second only \(p\).
All open shadows are nonempty and any
pair of them contains nondisjoint lieux,
so Mag L2 holds. They have the same
event structure and different refinement
and lieu data. The program checks Mag
1--4, divisibility and Mag L2 in both.
\end{proof}


\section{B.4. A faithful moderated o-minimal atelier}

\begin{definition}[Proved under stated hypotheses: the cell model]\label{SA21-BM-cell-definition}
Fix an o-minimal expansion \(\mathcal R\) of a real closed field \(R\),
an integer \(n\geq0\), and a definable set \(X\subseteq R^n\).
Definability always allows parameters from \(R\).
A \emph{generalized cell} is a nonempty definable set definably
homeomorphic to \((0,1)^d\) for \(d\geq0\), with \((0,1)^0\)
a singleton. No closed-ball condition is imposed on its closure.

A figure \(F\) is a finite partition of a closed definable
subset \(S\subseteq X\) into generalized cells, with frontier condition
\[
 D\cap\overline C^{\,X}\ne\varnothing
 \quad\Longrightarrow\quad D\subseteq\overline C^{\,X}
 \qquad(C,D\text{ cells of }F).
\]
The empty figure is included. Identify \(F\) with the finite
family of closed strata \(\{\overline C^{\,X}:C\in F\}\)
and write \(|F|=S\). For distinct cells \(D\subseteq\overline C^{\,X}\),
the frontier-dimension theorem
\citep[Theorem 3.22]{Coste1999} gives \(\dim D<\dim C\).
Thus their closures are distinct. The complement of \(C\)
in its closure is the finite union of the closed strata of
the lower cells, so \(C\) is relatively open and dense there.
The cell of a closed stratum \(A\) is
\[
 C_A=A\setminus\bigcup_{\substack{B\in F\\B\subsetneq A}}B.
\]
A multistrate \(m\) is the elementary figure of all strata
below one \(A\); put \(|m|=A\) and \(m^\circ=C_A\).
Thus it remembers the entire boundary decomposition.

Define \(G\leq F\) by downward-closed inclusion of closed-stratum
families, and define
\[
\begin{gathered}
 G\ll F\ \Longleftrightarrow\
 \text{each open cell of }G\text{ lies in an open cell of }F,\\
 G\preceq F\ \Longleftrightarrow\ G\ll F,\ |G|=|F|.
\end{gathered}
\]
On multistrates set \(m\perp^\circ n\) if
\(m^\circ\cap n^\circ=\varnothing\). Denote the resulting
structure by \(\mathcal A_{\mathrm{cell}}(X)\).
All closures are in the fixed ambient set \(X\), including
when \(X\) is unbounded or not closed in \(R^n\).
\end{definition}

\begin{lemma}[Proved: carriers, compatibility and restrictions]\label{SA21-BM-cell-calculus}
The relations \(\leq,\ll\) are partial orders with
\(\leq\subseteq\ll\). Refinement has a unique monotone carrier
\[
 c_{G,F}:\operatorname{Str}(G)\longrightarrow\operatorname{Str}(F),
 \qquad C_B\subseteq C_{c_{G,F}(B)}.
\]
For \(m\ll n\), the least substratum of \(n\) refined by
\(m\) is its top-cell carrier. Consequently
\[
 m\ll^\circ n\quad\Longleftrightarrow\quad
 m\ll n\text{ and }m^\circ\subseteq n^\circ.
\]
Figures are compatible exactly when, after uniting their
elementary subfigures, all distinct multistrates have disjoint
top cells; the join is then that union. They are disjoint
exactly when their supports are disjoint. For \(G\ll F\),
\(H\leq F\), induced refinement is
\[
 G|H=\{B\in G:c_{G,F}(B)\in H\},
 \qquad |G|H|=|G|\cap|H|.
\]
\end{lemma}

\begin{proof}
Reflexivity and transitivity follow from containment.
If \(F\ll G\ll F\), an \(F\)-cell \(C\) lies in a
\(G\)-cell \(D\), which lies in an \(F\)-cell \(C'\).
Nonempty intersection forces \(C=C'\), hence \(C=D\).
Repeating proves equality of figures. A downward-closed
subfamily retains its cells, so incidence implies refinement.
Carriers are unique since target cells are disjoint.
If \(B'\subseteq B\), then
\[
 C_{B'}\subseteq B'\subseteq B
 \subseteq\overline{C_{c(B)}}^{\,X}=c(B).
\]
The target cell \(C_{c(B')}\) meets \(c(B)\); frontier
compatibility gives \(c(B')\subseteq c(B)\). All carriers
of an elementary figure therefore lie below its top-cell
carrier, which must occur in any target subfigure containing
the refinement. This proves the carrier assertions.

For compatibility take the union of the proposed cells,
identifying equal elementary figures. Each original closed
stratum remains the union of its original downward family.
Frontier compatibility persists. No two distinct resulting
strata have equal support, since their disjoint top cells
would both be dense relatively open subsets of it.
Thus the union is a finite closed-support figure.
Conversely any common upper figure supplies precisely
these disjoint cells.
Disjoint supports permit their disjoint union as a figure,
with empty common incidence part. Conversely compatible
figures partition their supports by cells of their join;
a point in both lies in a common cell, and hence in a
nonempty common elementary subfigure.
Finally the inverse image of a downward family under
the carrier is downward. A point of \(|G|\) is in \(|H|\)
exactly when its carrier cell belongs to \(H\). This proves
the restriction formulas and their transitivity.
\end{proof}

\begin{theorem}[Proved: the complete axiom truth table]\label{SA21-BM-cell-axioms}
The chronological meanings of labels are fixed in Part B.1.
Positive entries hold for every definable \(X\); negative
entries mean failure for some such \(X\), proved below.
\begin{center}
\begin{tabular}{p{0.19\linewidth}p{0.57\linewidth}p{0.13\linewidth}}
\hline
System & Axioms & Value\\ \hline
156-4, Cont & Cont 1--5, Fig 1--4, finite geometric dimension,
finite principal ideals, Cond 7--8 & True\\
156-4, C & C1--C3, bounded-join primality, C5, C6,
C7\(_0\), canonical extension, C8, C8\(_0\) & True\\
156-4, C & Global incidence C4 and C4 bis & False\\
156-6 & At 1--8, At 9, At \(9'\), At 10--13,
At \(11'\), At 11 bis & True\\
156-6 & At \(9_\Sigma\) for all finite posets & True\\
156-6 & At \(9''\), or the empty/principal restriction
in \(9_\Sigma\) & False\\
156-6 & At 14, geometric At 14 bis, At 14 ter & True\\
156-6 & Literal all-target reflection in At 14 bis & False\\
156-7 & At C, At D, At D\(_0\), At CD, At CD2,
At CD2\( '\), At L1--L3, CL(spac), CL(pol), At supp & True\\
156-7 & At pol 1, At simp 1 & False\\
156-8 & Mag 1--4, strong-refinement formulation,
weaker primed clauses, Mog L1--L2, divisible-magasin clauses & True\\
156-8 & Atlas At 1--5, At Fil & True\\
156-9 & Renamed At 1--4, At spat 1--3 & True\\ \hline
\end{tabular}
\end{center}
Geometric At 14 bis means the mixed incidence condition
between a new boundary stratum and a stratum outside
that boundary. The literal all-target reading of its
uncertain transcription is treated separately below.
\end{theorem}

\begin{proof}
We verify mathematical contents, so changing labels
between chronological systems does not hide assumptions.

\emph{Multistrates.}
The orders and carriers give Mag 1--2.
Nonempty top cells make \(\perp^\circ\) symmetric and
irreflexive. Containment of top cells makes it hereditary
under strong refinement, proving Mag 3.
Distinct strata under one top stratum have disjoint
residual cells, proving the full all-pairs Mag 4.
Strong refinement is transitive since both refinement
and top-cell containment are transitive.
These facts give the corrected strong-refinement
formulation, its weakenings and renamed At 1--4.

\emph{Figures.}
The ideal below \(F\) is the downset lattice of its
finite stratum poset. Bounded joins are unions,
irreducibles are exactly elementary figures,
every figure is their join, and these irreducibles
are prime for bounded joins. The lemma proves
pairwise compatibility detection. These give
C1--C3 and At 1--3. Refinement is detected on
elementary figures, giving At 4. If \(G\leq F\),
\(m\leq F\), and \(m\ll G\), the top cell of \(m\)
is a cell of \(F\) contained in a cell of \(G\).
They are equal, so \(m\leq G\). More generally,
if compatible multistrates \(m,n\) satisfy \(m\ll n\),
put them in one common figure; the same cell
comparison gives \(m\leq n\). This is At 6 of
156-6. Least carriers give At 5, and the already
proved transitivity of strong refinement gives
At 7. Internal disjointness on figures
is defined only on elementary figures; this and
the elementary compatibility criterion give
both clauses of atlas At 5 in 156-8, scan 21.

\emph{Lieux and disjunction.}
Every singleton exists since parameters are allowed.
They are exactly the refinement-minimal nonempty
multistrates: every nonsingleton one has a proper
singleton refinement. Full and open shadows are
\(|m|\) and \(m^\circ\), respectively.
Every open shadow has a lieu; distinct lieux
are disjoint. The identity
\[
 m^\circ\cap n^\circ=\varnothing
 \quad\Longleftrightarrow\quad
 (\forall p\in m^\circ)(\forall q\in n^\circ)\ p\ne q
\]
and its whole-support analogue prove Mag L1--L2,
the displayed divisible-magasin conditions,
At 8--10, At L1--L3 and the lieu detection
clauses of the Cont systems.
Support inclusion makes disjointness hereditary
under refinement.

\emph{Compatibility and gluing.}
For figures \(F,G\), intersect their sets of elementary
subfigures to obtain their greatest common incidence
subfigure \(K\). Then
\[
 F\bowtie G\quad\Longleftrightarrow\quad
 (|F|\setminus|K|)\cap(|G|\setminus|K|)=\varnothing.
\]
The forward direction uses their common partition.
For the reverse, outside cells are disjoint by
the formula, inside cells are identical, and
outside cells cannot meet an inside cell.
The union criterion applies. This proves
CL(spac), At \(9'\), and At \(9_\Sigma\)
with all finite common posets allowed.

Suppose \(F\bowtie G\), \(F'\ll F\), \(G'\ll G\),
and \(F'|K\bowtie G'|K\). Any intersection
of a refining \(F'\)-cell and \(G'\)-cell
lies in \(|F|\cap|G|=|K|\). A refining cell
meeting \(|K|\) has its carrier in \(K\),
so belongs to the appropriate restriction.
All possible intersections are thus governed
by compatible restrictions. The union
criterion proves \(F'\bowtie G'\):
this is At 12 and At C.
The same restriction identity with
empty intersection proves At D and D\(_0\).
For absolute At CD2, suppose all refinements
of \(F,G\) whose restrictions to \(K\) are empty
are pairwise disjoint. Singleton points outside
\(|K|\) are among these refinements. Hence
\(|F|\setminus|K|\) and \(|G|\setminus|K|\)
are disjoint, and the displayed formula
gives \(F\bowtie G\). This proves At CD2;
the combined At CD formulation follows
from it and the relative criterion.
CL(pol) and At CD2\( '\) add a hypothesis
that \(K\) be empty or elementary to a
sufficient criterion, so follow from
CL(spac) and At CD2. They differ from
the biconditional At \(9''\).

Refinements agreeing on overlaps glue
by union, with precisely the prescribed
restrictions, proving At 11 and At \(11'\).
For subdivisions the supports cover the
original support, proving C8.
Any incidence cover of a finite figure
has a finite subcover. The data on
any further cover member agree with
the finite gluing on its intersections
with this subcover, which cover it.
Thus the arbitrary-cover version holds.
For At 11 bis, a directed family
with an existing incidence supremum
on one stratum is stationary there:
the supremum has finitely many
strata, included at one directed index.
Finitely many original strata give
one global stabilization index.
The eventual figure is the required
supremum. Starting instead with
stratumwise stationarity gives At Fil.

\emph{Supports and induction.}
If \(F'\ll F\) and some figure is disjoint
from \(F'\) but not \(F\), choose
\(p\in|F|\setminus|F'|\). Its singleton
refines \(F\), is disjoint from \(F'\),
and is not disjoint from \(F\).
This proves At supp, At 13 and
Cond 7 in their exterior-lieu forms.
It also proves that, for a refinement,
equal annihilators are equivalent to
equal supports, giving C5.
The restriction formula gives C6,
including uniqueness and commutation
with bounded joins. A refinement
missing a target point is strictly
enlarged by adjoining that singleton;
with equal supports no enlargement
is possible. Hence subdivisions are
exactly incidence-maximal refinements.
Antisymmetry gives C8\(_0\).
Extensions are proved separately.

\emph{Cont and dimension.}
The empty figure and points exist.
Admissibility, downward subfamilies,
compatibility detection and finite
compatible unions were proved in
the lemma. These verify the
preliminary and final Cont clauses.
Principal stratum ideals are finite.
Each strict frontier inclusion
lowers dimension by
\citet[Theorem 3.22]{Coste1999}.
For \(m'\ll m\), every chain
below \(m'\) therefore has length
at most \(\dim|m|\leq n\), since
\(|m'|\subseteq|m|\). This bounds
geometric dimension uniformly over
all refinements, as Cont 3 requires.
Empty figures, one-stratum figures
and dimension zero are included.
We do not assert that an unbounded
infinite union is a figure.
Spatiality and negative entries
are proved below.
\end{proof}


\begin{theorem}[Proved: faithfulness and the full powerset support lattice]\label{SA21-BM-cell-support}
The cell atelier is faithful and spatial. Its complete
support ortholattice for primitive internal disjointness is
\[
 \Sigma_M\cong\mathcal P(X).
\]
The open support of \(m\) corresponds to \(m^\circ\),
and its whole support to \(|m|\). This is the full
powerset, not only the definable subsets.
\end{theorem}

\begin{proof}
For arbitrary \(Q\subseteq X\), put
\(S_Q=\{m\in M:m^\circ\subseteq Q\}\).
For any \(\mathcal U\subseteq M\),
\[
 \mathcal U^\perp
 =S_{X\setminus\bigcup_{m\in\mathcal U}m^\circ}.
\]
Every point of \(Q\) supplies its singleton
multistrate in \(S_Q\), so
\[
 \bigcup_{m\in S_Q}m^\circ=Q,\qquad
 S_Q^\perp=S_{X\setminus Q},\qquad
 S_Q^{\perp\perp}=S_Q.
\]
Every closed polar has this form.
Thus \(Q\mapsto S_Q\) is a bijection
preserving inclusion and complement,
hence all complete Boolean operations.
This includes \(X=\varnothing\).
Nondefinable \(Q\) are permitted because
polarity allows arbitrary multistrate
families.

Distinct lieux are disjoint, and
primitive disjointness is exactly
disjointness of their open shadows,
proving faithfulness.
For spatiality, the ordered deployment
\(\{S_{n^\circ}:n\leq m\}\)
recovers each cell \(n^\circ\), since
\(Q\mapsto S_Q\) is injective.
Downward unions recover each closed
support, hence the whole elementary
figure \(m\). This proves
At spat 3 of 156-9, scan 50.
The incidence order and symmetric
irreflexive relation give At spat 1.
Distinct incident strata have
disjoint cells, giving At spat 2.
\end{proof}

\begin{theorem}[Proved: coarsest and elementary extensions]\label{SA21-BM-cell-extension}
If \(G\leq F\), \(G'\preceq G\), the closed-stratum family
\[
 E=G'\cup(F\setminus G)
\]
is a figure with \(E\preceq F\), \(E|G=G'\).
Every other extension \(H\preceq F\) inducing \(G'\)
refines \(E\), so \(E\) is the unique coarsest extension.
Its carrier is an order isomorphism outside the
prescribed boundary and has the required mixed
incidences. Every extension with this complement
property is equal to \(E\).
For elementary \(m\), a subdivision of \(\partial m\)
extends uniquely to an elementary \(m'\preceq m\)
with that boundary. Thus At 14, geometric
At 14 bis and At 14 ter hold.
\end{theorem}

\begin{proof}
No \(A\in F\setminus G\) lies in \(|G|\),
since its open cell is disjoint from \(|G|\).
Hence it lies below no stratum of \(G'\).
The residual cell of every new stratum
is therefore unchanged. For an old
\(A\in F\setminus G\), the old boundary
portion in \(G\) is \(A\cap|G|\).
Its induced subdivision by \(G'\)
covers exactly the same set.
Other old boundary strata remain,
so the residual cell of \(A\)
also remains unchanged.
The resulting cells are the
cells of \(G'\) and the old
cells outside \(G\).
They partition \(|F|\), retain
their stated closures and obey
the frontier condition by this
boundary calculation. They
refine \(F\), proving the claims
about \(E\).

For another extension \(H\),
an \(H\)-cell with carrier
in \(G\) belongs to \(H|G=G'\);
every other cell lies in an
old unchanged cell outside
\(G\). Thus \(H\ll E\).
Antisymmetry proves uniqueness
of a coarsest extension.
On the outside strata the
carrier of \(E\) is the
identity. For \(B\in G'\)
and \(A\in F\setminus G\),
\[
 B\subseteq A\quad\Longleftrightarrow\quad c(B)\subseteq A.
\]
One implication uses \(B\subseteq c(B)\).
For the other, the old carrier
cell of \(B\) meets \(A\);
frontier compatibility puts
its closure in \(A\).
These are the mixed incidences.
If another extension has a
bijective outside carrier,
only one cell covers each
old cell outside \(G\).
Equal support forces that
cell to be the whole old
cell, with its old closure.
Its inside figure is specified,
so the extension is \(E\).

For elementary \(m\), use
\(G=\partial m\).
The construction retains
the one top closed stratum
\(|m|\), hence is elementary
with top cell
\(|m|\setminus|\partial m|=m^\circ\).
Any other elementary extension
with the given boundary has
the same support and boundary
support, hence the same top
cell and all the same strata.
This proves uniqueness.
No triangulation is needed;
allowing general cell closures
rather than simplices is
what makes the construction
possible.
\end{proof}

\begin{counterexample}[Refuted: literal all-target At 14 bis]\label{SA21-BM-14bis-reading}
The literal mixed-order clause of 156-6, scan 88,
as transcribed and visually checked in Part B.1,
is false in the cell model if its second variable
is allowed to range over all of \(\widetilde F'\).
The corrected clause restricts that variable
to \(\widetilde F'\setminus\widetilde G'\).
\end{counterexample}

\begin{proof}
Take \(G=F\) the interval figure on \([0,1]\),
with the two endpoint cells and the open interval.
Subdivide at \(1/2\) to get \(G'\).
An extension inducing \(G'\) must be \(F'=G'\),
since \(G=F\). The two distinct open half-interval
cells have the same old interval carrier.
Literal order reflection for all pairs
of boundary strata would make each
of their closed strata lie below the other,
which is false. The five-stratum carrier
poset is completely specified and checked
by the program. The corrected mixed
condition, with the second stratum
outside \(G'\), is proved by the
previous theorem, as are the other
cartesian and complement conditions.
Thus the literal reading and its
geometric correction have separate
truth values.
\end{proof}

\begin{theorem}[Known input and proved consequence: moderation]\label{SA21-BM-cell-moderation}
Every finite collection of subdivisions of a cell figure
has a common subdivision. Every closed definable
\(S\subseteq X\) has a cell figure, and any finite
collection of its finite cell stratifications
has a common refinement. In particular the
atelier is modérée in the sense of 156-4, scan 70.
\end{theorem}

\begin{proof}
The input is compatible definable triangulation:
if \(A\subseteq R^n\) is closed, bounded and
definable and \(B_1,\ldots,B_r\subseteq A\)
are definable, there are a finite simplicial
complex \(K\) and a definable homeomorphism
\(h:|K|\to A\) such that every \(B_i\)
is a union of images of open simplices.
We use the fully opened statement of
\citet[Theorem 4.4, \S4.3, p.~42;
proof pp.~42--44]{Coste1999}, over an
arbitrary fixed o-minimal expansion
of a real closed field.
The standard book locator is
\citet[Chapter 8, \S2, Theorem 2.9]{vandenDries1998};
its theorem page was not obtained
in full and is not the evidential
basis of this proof.

Use the definable homeomorphism
\[
 b:R^n\longrightarrow(-1,1)^n,\qquad
 b(t)_i=\frac{t_i}{1+|t_i|},
\]
whose inverse is \(s_i/(1-|s_i|)\).
Let \(A=\overline{b(X)}^{\,R^n}\).
Triangulate it compatibly with
\(b(X)\), \(b(S)\), and the images
of all cells in the given finite
collection of stratifications.
For every open simplex \(\sigma\)
whose image lies in \(b(S)\), set
\[
 C_\sigma=b^{-1}(h(\sigma)).
\]
It is definably homeomorphic to
an open simplex, hence to
\((0,1)^{\dim\sigma}\) by a
semialgebraic homeomorphism.
The selected cells partition \(S\).

The closure of \(h(\sigma)\)
in \(b(X)\) is
\(h(\overline\sigma)\cap b(X)\).
An open face meeting \(b(X)\)
lies wholly in it by compatibility.
Since \(b(S)\) is closed in
\(b(X)\), every such face lies
in \(b(S)\). Thus each cell
closure in \(X\) is a union
of selected full cells, with
the frontier condition.
The closed strata form a
figure \(H\) in our exact
model. Compatibility with
all old cells makes every
new cell lie in one old
cell. It follows that
\(H\) refines all the
given figures and has
their support \(S\).
This handles unbounded
or nonclosed \(X\), and
\(S\) closed only in \(X\).
If \(S=\varnothing\), use
the empty figure.

No common linear subdivision
of fixed triangulating maps
is needed. Compatibility
with finitely many definable
cells is sufficient; no
Hauptvermutung or uniqueness
theorem is used.
\end{proof}

\begin{proposition}[Proved: the correct connectedness statement]\label{SA21-BM-cell-connected}
The undirected refinement components of \((M,\ll)\)
are exactly the definably connected components
of \(X\): a multistrate is assigned the unique
component containing its support, and two
multistrates are in the same refinement
component exactly when those assignments
agree. Different refinement components
are compatible. For an individual finite
figure, its incidence components are
exactly the definably connected components
of its support.
\end{proposition}

\begin{proof}
Generalized cells are definably connected,
since open boxes are definably path
connected. A closure of a definably
connected set is definably connected:
a definable clopen partition would
restrict to the dense set, and
density excludes an unused nonempty
open part. Thus \(|m|\) lies in
one component. Refinement implies
support inclusion, so the assigned
component is constant on refinement
edges. Each \(p\in|m|\) also
gives \(\{p\}\ll m\).

Conversely
\citet[Theorem 3.9, p.~28]{Coste1999}
gives finitely many definably
connected components, each definably
arcwise connected. Join \(p,q\)
in one component by a continuous
definable \(\gamma:[0,1]\to X\);
no injectivity is needed.
Its image \(P\) is bounded by
the extreme-value property
applied to its coordinates,
and closed by the same property
applied to distance from a
proposed missing closure point.
These follow from
\citet[Theorem 3.4, pp.~26--27]{Coste1999}.
The image is definably connected
as a continuous image of \([0,1]\).

Triangulate \(P\) compatibly
with its endpoints. The
resulting finite cell figure
in \(X\) has \(p,q\) as
vertices and connected
incidence graph: otherwise
an incidence component
and its complement would
be finite unions of closed
strata, giving a nontrivial
definable clopen partition.
Incidence implies refinement,
so this graph connects the
points and then the original
multistrates. Distinct
components have disjoint
supports, hence compatibility.

For an individual figure,
incidence components are
definably clopen by the
same finite-union argument.
They are connected because
each closed stratum is
connected and incident
strata intersect. Their
connected finite incidence
graph gives connected union,
successively along a spanning
tree. These are therefore
exactly its support components.
\end{proof}

\begin{proposition}[Refuted stronger assertions and proved scope limits]\label{SA21-BM-cell-limits}
Global incidence C4 and C4 bis, the polyhedral
restrictions and an infinite-family strengthening
of moderation fail in general. Allowing arbitrary
nonclosed supports with closures taken only
relative to themselves also destroys the
disjoint-support compatibility axiom.
\end{proposition}

\begin{proof}
For C4 let \(X=R\). The one-cell
figure on \(R\) has no proper
incidence substratum. It cannot
occur properly below another
elementary figure since its
support is already all of \(R\).
It is isolated in the global
incidence poset. The singleton
at \(0\) lies in a different
component and is incompatible
with it, since the open cells
intersect. C4 fails, and so
does its strengthening C4 bis.

For the polyhedral restrictions,
partition the semialgebraic unit
circle into two points \(p,q\)
and two complementary open arcs.
The arc closures intersect in
the two-point figure, neither
empty nor elementary. This
refutes At pol 1 and the
biconditional At \(9''\).
It does not refute CL(pol)
or At CD2\( '\), which
are sufficient criteria.
Partitioning
the circle instead into
one point and its complement
gives a top stratum with
exactly one immediate predecessor,
refuting At simp 1.
The cells are generalized
one-cells: stereographic
projection identifies a
punctured circle with \(R\),
then with \((0,1)\); the
two arcs are semialgebraically
homeomorphic to intervals.
The finite incidence tests
are in the program; these
geometric realizations put
the checked posets in the
specified model.

For arbitrary nonclosed supports,
take the one-cell partition
of \((0,1)\subset R\) and the
singleton at \(1\).
Their supports are disjoint,
but they cannot both be
closed subfigures of a
common frontier stratification:
\((0,1)\) is not closed in
\((0,1]\), so its cell
acquires the new boundary.
This is the obstruction
removed by fixed ambient
closures and closed supports.

Finally subdivide the standard
interval figure at every
\(t\in(0,1)\), one subdivision
for each \(t\). A common
refinement with the same
support would have every
singleton \(t\) as a cell,
since it must cover each
point carrier. It would
be infinite, contrary to
the figure definition.
Moderation only quantifies
over pairs. Disconnected
supports cause no problem:
their finitely many definable
components can be stratified
separately and united.
Disconnected sets are
not called individual
generalized cells.
\end{proof}
